\begin{figure}[htbp]
\centering
\begin{tikzpicture}[fa diagram]
  \draw[fa axis] (-.6,0)--(7,0) node[right] {$M$};
  \coordinate (O) at (0,0); \coordinate (P) at (3,0); \coordinate (X) at (3,2.6); \coordinate (Z) at (6,0);
  \draw[->,fa main] (O)--(X);
  \node[above] at (X) {$x$};
  \draw[->,fa main] (O)--(P);
  \draw[fa accent] (P)--(X);
  \node[rotate=90,anchor=south] at (3.12,1.3) {$x-P_Mx$};
  \draw[fa main,dashed] (X)--(Z) node[midway,right] {$x-z$};
  \draw[fa main] (3,.25)--(3.25,.25)--(3.25,0);
  \node[below] at (P) {$P_Mx$}; \node[below] at (Z) {$z$};
  \node[fa note] at (3.2,-1) {$\|x-z\|^2=\|x-P_Mx\|^2+\|P_Mx-z\|^2$};
\end{tikzpicture}
\caption{正交投影的勾股分解。$M$ 闭且 Hilbert 空间完备，才保证垂足存在；对任意 $z\in M$，第二项非负，所以投影是唯一最近点。图画出二维实截面，复 Hilbert 空间使用同一范数恒等式。}
\label{fig:y3-projection}
\end{figure}
